\documentclass[10pt,a4paper]{amsart}
\usepackage[margin=19mm]{geometry}
\usepackage{amsmath,amssymb,mathtools}
\usepackage[scr=rsfs,cal=euler]{mathalfa}
\usepackage{xfrac}
\usepackage[unicode,hidelinks,bookmarks=false]{hyperref}
\newcommand{\ie}{i.e.\ }
\newcommand{\cf}{cf.\ }
\newcommand{\resp}{respectively\ }
\newcommand{\Subsection}{\subsection}
\providecommand{\nobreakdash}{\nobreak\hbox{-}}
\setlength{\emergencystretch}{2em}
\multlinegap0pt
\usepackage{lmodern}
\usepackage{booktabs}
\usepackage{algorithm}
\usepackage{algorithmic}
\usepackage{bm}
\usepackage{mathtools}
\usepackage{dsfont}
\newcommand{\ind}{\mathds{1}}
\newcommand{\defn}{\coloneqq}
\newcommand{\defnrev}{\eqqcolon}
\newcommand{\var}{\mathsf{Var}}
\newcommand{\cov}{\mathsf{Cov}}
\newcommand{\cX}{\overline{X}}
\newcommand{\cS}{\mathcal{S}}
\newcommand{\cA}{\mathcal{A}}
\newcommand{\cB}{\mathcal{B}}
\newcommand{\cC}{\mathcal{C}}
\newcommand{\cD}{\mathcal{D}}
\newcommand{\cE}{\mathcal{E}}
\newcommand{\cF}{\mathcal{F}}
\newcommand{\cG}{\mathcal{G}}
\newcommand{\cH}{\mathcal{H}}
\newcommand{\cI}{\mathcal{I}}
\newcommand{\cJ}{\mathcal{J}}
\newcommand{\cK}{\mathcal{K}}
\newcommand{\cT}{\mathcal{T}}
\newcommand{\cP}{\mathcal{P}}
\newcommand{\cN}{\mathcal{N}}
\newcommand{\cM}{\mathcal{M}}
\newcommand{\bO}{\mathbb{O}}
\newcommand{\bT}{\mathbb{T}}
\newcommand{\bE}{\mathbb{E}}
\newcommand{\bP}{\mathbb{P}}
\newcommand{\bR}{\mathbb{R}}
\newcommand{\bN}{\mathbb{N}}
\newcommand{\bS}{\mathbb{S}}
\newcommand{\bZ}{\mathbb{Z}}
\newcommand{\bX}{\mathbb{X}}
\newcommand{\KL}{\mathsf{KL}}
\newcommand{\diff}{\,\mathrm{d}}
\newcommand{\row}{\mathsf{row}}
\newcommand{\col}{\mathsf{col}}
\newcommand{\numpf}[2]{\overset{(\mathrm{#1})}{#2}}
\newcommand{\ol}{\overline}
\newcommand{\wh}{\widehat}
\newcommand{\wt}{\widetilde}
\newcommand{\clip}{\mathsf{clip}}
\newcommand{\data}{\mathsf{data}}
\newcommand{\sde}{\mathsf{sde}}
\newcommand{\ode}{\mathsf{ode}}
\newcommand{\deter}{\mathsf{det}}
\newcommand{\tr}{\mathsf{tr}}
\newcommand{\veps}{\varepsilon}
\newcommand{\vphi}{\varphi}
\newcommand{\TV}{\mathsf{TV}}
\newcommand{\setc}{\mathrm{c}}
\newcommand{\md}{\mathsf{mid}}
\newcommand{\acc}{\mathsf{acc}}
\newcommand{\score}{\mathsf{score}}
\newcommand{\disteq}{\overset{\mathrm d}=}
\newcommand{\aux}{\mathsf{aux}}
\newcommand{\jcb}{\mathsf{jcb}}
\newcommand{\thr}{\mathsf{clip}}
\newcommand{\gd}{\nabla}
\newcommand{\ul}{\underline}
\newcommand{\polylog}{\mathsf{poly}\log}
\newcommand{\sm}{\mathsf{sm}}
\newcommand{\typ}{\mathsf{typ}}
\newcommand{\psd}{\succcurlyeq}
\newcommand{\es}{\mathsf{es}}
\newcommand{\proj}{\mathsf{proj}}
\newcommand{\cor}{\mathsf{cor}}
\newcommand{\pre}{\mathsf{for}}
\newcommand{\back}{\mathsf{bck}}
\newcommand{\supp}{\mathsf{supp}}
\newcommand{\low}{\mathsf{low}}
\newtheorem{lemma}{\bf Lemma} 
\newtheorem{proposition}{\bf Proposition}
\newtheorem{theorem}{\bf Theorem}
\newtheorem{corollary}{\bf Corollary} 
\newtheorem{claim}{\bf Claim}
\newtheorem{assumption}{\bf Assumption} 
\newtheorem{definition}{\bf Definition} 
\newtheorem{condition}{\bf Condition}
\newtheorem{property}{\bf Property} 
\newtheorem{example}{\bf Example}
\newtheorem{fact}{\bf Fact}
\newtheorem{remark}{Remark}
\begin{document}
\title[Diffusion Models Are Statistically Optimal for Learning Low-]{Diffusion Models Are Statistically Optimal for Learning Low-Dimensional Multi-Modal Distributions}
\author{Jingda Wu, Changxiao Cai}
\date{}
\maketitle
\noindent\emph{Source and licence.} Diffusion Models Are Statistically Optimal for Learning Low-Dimensional Multi-Modal Distributions. Version arXiv:2605.30153v1 (CC BY 4.0). This material is distributed under the Creative Commons Attribution 4.0 International licence, http://creativecommons.org/licenses/by/4.0/, as recorded in the arXiv OAI licence metadata for this exact version and shown on the abstract page https://arxiv.org/abs/2605.30153v1. Full rights notice: licenses/arxiv-2605.30153-CC-BY-4.0.txt.\par\smallskip
\noindent\textit{Editorial scope.} Counted unit: the original \texttt{lemma} environment \texttt{lem:tail bdd for subgaussian} at source lines 1396--1407 together with its complete proof at lines 1409--1427; the delivered document keeps source lines 1394--1428 so that every referenced label and every citation resolves inside it. Native editable source is the paper's own concatenated TeX; the AMS wrapper is editorial formatting only. No publisher class, upstream script or unrelated artwork is redistributed.
\par\smallskip\noindent\textit{Reference note.} This article ships no compiled bibliography (biblatex); the entries delivered here are composed from the author's own BibTeX (.bib) fields, with no field invented and no entry dropped.
\section{Auxiliary Lemmas}
% \subsection{Tail bound for subgaussian r.v}

\begin{lemma}[Tail bound for subgaussian random vectors]
        \label{lem:tail bdd for subgaussian}
        Let $\nu$ be a $\sigma$-subgaussian distribution in $\mathbb{R}$, i.e, 
    \begin{align}
        \sigma = \| X\|_{\psi_2} \defn \inf \Big\{ t>0: \mathbb{E}\exp \big(X^2/t^2\big) \leq 2\Big\}, \quad X\sim \nu\label{eq:subgaussian def}.
    \end{align}
    Then we have the following tail bound for $X \sim \nu$, any $r\geq 0$ and some absolute constant $c>0$, 
    \begin{align}
        \mathbb{P}[|X| \geq r] &\leq  2\exp\big(-\frac{cr^2}{\sigma^2}\big) \label{eq:subgaussian tail bdd} \\
        \mathbb{E}[|X|^2 \ind_{\{|X|\geq r\}}] & \leq 2 (r^2 + \sigma^2/c) \exp\big(-\frac{cr^2}{\sigma^2}\big) .\label{eq:subgaussian tail bdd of second moment}
    \end{align}
\end{lemma}

\begin{proof}
\begin{itemize}
    \item \textbf{Proof of (\ref{eq:subgaussian tail bdd}).}
    This is a classical result for sub-gaussian r.v. We can find the proof in Proposition 2.5.2 in \cite{Vershynin_2018}. 

\item \textbf{Proof of (\ref{eq:subgaussian tail bdd of second moment}).}
For any $r\geq0$, 
\begin{align*}
    \mathbb{E} \big[ |X|^2\ind_{\{|X|\geq r\}} \big] &= \int_{\mathbb{R}} x^2 \ind_{\{|x|\geq r\}} \cdot p(x) \!\diff x \\
    & = \int_{\mathbb{R}} \Big( \int_0^{|x|} 2z dz\Big) \ind_{\{|x|\geq r\}} p(x) \!\diff x \\
    &= 2\int_{\mathbb{R}}  \int_0^{+\infty} z\cdot \ind_{\{|x|\geq z\}}  \ind_{\{|x|\geq r\}} p(x) \!\diff z \!\diff x \quad (\text{Tonelli's Theorem}) \\
    &= 2\int_0^{+\infty} \Big( \int_{\mathbb{R}}\ind_{\{|x| \geq r\vee z\}}\cdot  p(x) \!\diff x\Big)z \!\diff z \quad (\text{Tonelli's Theorem})\\ 
    & \leq r^2 \mathbb{P}[|X| \geq r] + 2 \int_r^{+\infty} z \cdot \mathbb{P}[|X| \geq z] \!\diff z \\
    &\leq r^2 \mathbb{P}[|X| \geq r] + 4 \int_r^{\infty} z \cdot e^{-cz^2/\sigma^2}  \!\diff z \\
    &= 2 (r^2 + \sigma^2/c) \exp\big(-\frac{cr^2}{\sigma^2}\big) .
\end{align*}
\end{itemize}

\end{proof}


\begin{thebibliography}{99}
\bibitem[]{Vershynin_2018} Vershynin, Roman. High-dimensional probability. University of California, Irvine 10 (2020) 31
\end{thebibliography}
\end{document}
